\section{Welfare and sale terms}\label{sec:welfare}

\subsection{Access to prices}\label{sec:access}

Access to prices improves acquisition outcomes at a fixed level of competition. Hold the incumbent at \(F_1\).
Compare the equilibrium of Proposition~\ref{prop:ccc} with that of an economy in which the challenger does not see the
price. Under the hypotheses of Proposition~\ref{prop:ccc}, both economies have unique full orders. Price access raises
expected target proceeds and acquisition surplus net of preparation costs (Proposition~\ref{prop:access}). Acquisition
surplus excludes transfers between bidders, shareholders and traders.

The two economies generate the same orders. The trading bound of Lemma~\ref{lem:trading} forces full orders even
when the challenger cannot use the price and entry stays at \(\rho\). Trading costs therefore coincide. Price access
changes only which high-cost challengers prepare: it adds preparation after a purchase. Each added preparation
contributes its gross profit minus its cost, and also the sales that the reserve would otherwise block. A challenger
prepares only when its expected gross profit covers its cost, so each term is positive.

In the benchmark at \(r_1\), expected proceeds rise from \(\reg{p2_base_hidden_proceeds}\) to
\(\reg{p2_base_r1_proceeds}\), a gain of \(\reg{p2_access_proceeds_gain}\). Net surplus rises by
\(\reg{p2_access_surplus_gain}\). With research at the cap (Section~\ref{sec:research}), the gains are the same.

Target shareholders gain because a shareholder who holds through the sale receives \(V_T\), not the interim price.
The informed investor's expected gross profit, \(\reg{p2_access_gross_transfer}\) at \(r_1\), is a transfer from
liquidity traders. Suppose that one order unit is \(z\) shares per share of float. Relative to an economy without the
informed investor, initial shareholders gain in total at every \(z<\reg{p2_access_break_even}\), even if they bear the
whole transfer.

The gain needs the recruitment region of Proposition~\ref{prop:ccc}, in which high-cost preparation does not pay at
the prior. When it already pays at the prior, an informative price deters some preparation. Keep the strong benchmark
and set \(c_H=\reg{access_out_cH}\): price access then changes proceeds by \(\reg{p2_outside_proceeds_gain}\) and net
surplus by \(\reg{p2_outside_surplus_gain}\) (Appendix~\ref{app:welfare}). The information in the price, not its
level, separates the two economies. A deterministic external dividend of \(\reg{p2_access_proceeds_gain}\) in the
hidden-price economy equalizes the mean prices, and it leaves the orders and entry unchanged.

\subsection{The payment rule}\label{sec:payment}

The effect of competition on the target-payoff spread depends on the sale payment. Proposition~\ref{prop:payment}
isolates this dependence in a verifiable-value institution with zero reserve. After diligence, a sale to the next
buyer at its value is an enforceable fallback. The seller and the winner divide the surplus above the fallback by Nash
bargaining, with seller weight \(\omega\in[0,1)\). A stronger incumbent weakly lowers the challenger's profit, and
its effect on the spread weakly has the sign of \(1-2\omega\). This effect is strict when \(\E[(R-\ell)_+]\) changes
strictly and \(\omega\ne1/2\); at \(\omega=1/2\) the spread does not depend on the incumbent.

To see why, take \(R>\ell\). A high-class challenger wins and pays \((1-\omega)R+\omega h\), whose sensitivity to
\(R\) is \(1-\omega\). A low-class challenger loses, and the incumbent pays \((1-\omega)\ell+\omega R\), whose
sensitivity is \(\omega\). Competition widens the difference exactly when the first sensitivity exceeds the second.

In the benchmark, the spread under this rule changes by \(\reg{pay_gap_coef}\,(1-2\omega)\) from \(r_0\) to \(r_1\).
The second-price auction is the case \(\omega=0\): it puts the whole weight on the runner-up's value.
\citet{BulowKlemperer1996} compare an auction with one more bidder and an optimal negotiation by expected revenue;
here the comparison concerns how competition moves the spread. The trading and preparation equilibrium under the
bargaining rule is open.

\subsection{Reserves}\label{sec:reserves}

A reserve affects both the payment required to win and the information in the target's shares. Once the reserve can
exclude bidders, three events separate. \emph{Preparation} is the event that the challenger pays its preparation
cost. A \emph{sale} occurs when at least one bidder submits an admissible bid. \emph{Two admissible bidders} meet
when both the incumbent and a prepared challenger meet the reserve. In the benchmark of Section~\ref{sec:benchmark},
the reserve lies below both challenger values, so every prepared challenger is admissible.

A reserve that excludes the low class can widen the class difference in target proceeds, because the low class then
adds nothing to proceeds. Informed trade can then pay against both incumbents, and the price is informative at both
strengths. Proposition~\ref{prop:reserves} gives the consequence. If the top posterior recruits at both strengths,
entry is flat in incumbent strength. Strict deterrence needs the stronger incumbent to push gross profit at the top
posterior below the high cost.

\begin{table}[tbp]
\centering\small\singlespacing
\caption{Reserves, preparation and sale in the benchmark.}\label{tab:reserves}
\begin{tabular}{@{}llcccc@{}}
\toprule
Reserve & Strength & Preparation & Sale & Two admissible bidders & Expected proceeds\\
\midrule
\(\reg{prim_p}\) & \(r_0\) & \reg{p2_reserve_points_base_r0_entry} & \reg{p2_reserve_points_base_r0_sale} & \reg{p2_reserve_points_base_r0_two} & \reg{p2_reserve_points_base_r0_proceeds}\\
 & \(r_1\) & \reg{p2_reserve_points_base_r1_entry} & \reg{p2_reserve_points_base_r1_sale} & \reg{p2_reserve_points_base_r1_two} & \reg{p2_reserve_points_base_r1_proceeds}\\
\(\reg{rsv_points_excl}\) & \(r_0\) & \reg{p2_reserve_points_exclude_r0_entry} & \reg{p2_reserve_points_exclude_r0_sale} & \reg{p2_reserve_points_exclude_r0_two} & \reg{p2_reserve_points_exclude_r0_proceeds}\\
 & \(r_1\) & \reg{p2_reserve_points_exclude_r1_entry} & \reg{p2_reserve_points_exclude_r1_sale} & \reg{p2_reserve_points_exclude_r1_two} & \reg{p2_reserve_points_exclude_r1_proceeds}\\
\(\reg{rsv_deter}\) & \(r_0\) & \reg{p2_reserve_points_deter_r0_entry} & \reg{p2_reserve_points_deter_r0_sale} & \reg{p2_reserve_points_deter_r0_two} & \reg{p2_reserve_points_deter_r0_proceeds}\\
 & \(r_1\) & \reg{p2_reserve_points_deter_r1_entry} & \reg{p2_reserve_points_deter_r1_sale} & \reg{p2_reserve_points_deter_r1_two} & \reg{p2_reserve_points_deter_r1_proceeds}\\
\bottomrule
\end{tabular}
\par\medskip
\begin{minipage}{0.92\linewidth}\footnotesize
\textit{Notes.} Benchmark \(B\) with point-mass classes. At each reserve, full informed orders are the unique outcome
at both strengths, except at the reserve \(\reg{prim_p}\) and \(r_0\), where the investor does not trade. Values are
computer-assisted.
\end{minipage}
\end{table}

Table~\ref{tab:reserves} reports three reserves in the benchmark (computer-assisted). The reserve
\(\reg{rsv_points_excl}\) excludes the low class. Full orders are unique at both strengths, only a purchase recruits,
and entry is \(\reg{p2_reserve_points_exclude_r0_entry}\) at both \(r_0\) and \(r_1\). This reserve raises expected
proceeds at both strengths. Narrow uniform classes give the same pattern at the reserve \(\reg{rsv_bands_excl}\).

The reserve \(\reg{rsv_deter}\) restores strict deterrence: entry is \(\reg{p2_reserve_points_deter_r0_entry}\) at
\(r_0\) and \(\reg{p2_reserve_points_deter_r1_entry}\) at \(r_1\), and proceeds exceed their levels at the reserve
\(\reg{prim_p}\). This reserve excludes the low class and also the weak incumbent, since it exceeds \(r_0\). In the
benchmark, strict deterrence after the low class is excluded requires this second exclusion. With a reserve in
\((\ell,r_0]\), the top posterior recruits at both strengths, and entry is flat. Strict deterrence holds for reserves
between \(p_H\approx\reg{p2_reserve_ph}\) and \(h-c_H/M\approx\reg{p2_reserve_weak_top}\) (Appendix~\ref{app:welfare}).
The seller's choice of a reserve with the whole continuation retained is open.
